% Реферат по ГОСТ/методичке при необходимости подключают отдельно (не в начале основного текста).
% Заголовок и запись в оглавление добавьте в main.tex по требованиям кафедры.
\section*{Реферат}

Объектом исследования в работе является процесс оркестрации контейнерных приложений в распределенной инфраструктуре с использованием Docker Compose, центрального управляющего узла и набора удаленных агентных узлов. Предметом исследования выступают архитектурные, алгоритмические и программные решения, обеспечивающие публикацию, запуск, контроль состояния и динамическую TCP-маршрутизацию сервисов.

Целью работы является разработка и исследование программной платформы, обеспечивающей воспроизводимое развертывание Compose-проектов, централизованное управление агентами и проксирование внешних TCP-подключений к сервисам, запущенным на удаленных хостах.

В рамках работы решены следующие задачи: проведен анализ существующих подходов к оркестрации; сформулированы функциональные и нефункциональные требования к системе; разработана архитектура платформы Chronos; реализованы серверные и клиентские компоненты (Chronos.Master, Chronos.Agent, Chronos.Core, Chronos.Master.Web); реализованы резервное копирование томов в объектное хранилище по политикам Master, механизм генерации конфигурации HAProxy для TCP-маршрутов и отправка логов в Loki; проведены функциональные испытания и проанализированы ограничения.

Практическая значимость работы состоит в возможности применения платформы как инженерного контура для малых и средних команд, использующих Docker Compose в сценариях CI/CD, стендов предрелизного тестирования и операционного сопровождения.

Ключевые слова: оркестрация контейнеров, Docker Compose, распределенная система, агентная архитектура, TCP-проксирование, HAProxy, резервное копирование томов, .NET, React, Loki, Grafana.
